% =======================
% Section 3: Agent Patterns
% =======================
\section{Agent Patterns}

% ---- Slide: ReAct Framework ----
\begin{frame}{ReAct Framework}
\begin{columns}[T]
\begin{column}{0.52\textwidth}
\textbf{ReAct} (Reasoning + Acting)
\begin{itemize}\setlength\itemsep{0.2em}
  \item The canonical single-agent execution pattern
  \item Alternates between:
  \begin{itemize}
    \item \textbf{Thought} --- reason about the next step
    \item \textbf{Act} --- take an action
    \item \textbf{Observe} --- receive result, update beliefs with real-world feedback
    \item \textbf{Repeat} until goal is reached
  \end{itemize}
  \item The \textbf{Observe} step is what a static LLM cannot do - grounding in real-world feedback
\end{itemize}
\end{column}
\begin{column}{0.44\textwidth}
\centering
\begin{tikzpicture}[
  node distance=0.7cm,
  stepnode/.style={draw, circle, fill=blue!12, minimum size=1.1cm,
               align=center, font=\small},
  every node/.style={font=\small}
]
  \node[stepnode] (goal)   at (0, 3.2) {\textbf{Goal}};
  \node[stepnode, fill=orange!20] (thought) at (0, 2.2) {Thought};
  \node[stepnode, fill=green!15]  (act)     at (0, 1.2) {Act};
  \node[stepnode, fill=yellow!20] (obs)     at (0, 0.2) {Observe};
  \node[draw, rounded corners=3pt, fill=red!10,
        minimum width=1.8cm, minimum height=0.55cm,
        font=\small] (stop) at (0,-0.65) {Done};

  \draw[-{Stealth}] (goal)   -- (thought);
  \draw[-{Stealth}] (thought) -- (act);
  \draw[-{Stealth}] (act)    -- (obs);
  \draw[-{Stealth}, dashed]  (obs) -- node[right, font=\scriptsize]{or loop} (stop);
  \draw[-{Stealth}] (obs.east) to[out=0,in=0,looseness=2.2] (thought.east);
\end{tikzpicture}
\vspace{0.3em}

\scriptsize Other patterns: \textbf{ReWOO} (plan ahead, no observation loop), \textbf{Reflexion} (self-reflection on errors), \textbf{LATS} (tree search + self-reflection)
\end{column}
\end{columns}
\end{frame}


% ---- Slide: Orchestration Architecture ----
\begin{frame}{Orchestration: Single-Agent vs.\ Multi-Agent}
\begin{columns}[T]
\begin{column}{0.50\textwidth}

\begin{block}{Multi-Agent System}
\begin{itemize}
  \item \textbf{Orchestrator agent}: decomposes the goal, delegates to subagents, synthesizes results
  \item \textbf{Subagents}: specialized, each running its own internal loop
  \item Enables parallelism and role separation
  \item Overhead: coordination, communication, trust between agents
\end{itemize}
\end{block}
\end{column}
\begin{column}{0.46\textwidth}
\begin{tikzpicture}[
  box/.style={draw, rounded corners=3pt, minimum width=2.8cm,
              minimum height=0.6cm, align=center, font=\small},
  sbox/.style={draw, rounded corners=3pt, minimum width=2.0cm,
               minimum height=0.55cm, align=center, font=\scriptsize},
  every node/.style={font=\small}
]

  % Multi-agent
  \node[box, fill=red!10] (orch) at (0, 2.3) {Orchestrator Agent};
  \node[sbox, fill=orange!15] (rl0) at (0, 1.55) {ReAct loop};
  \draw[-{Stealth}] (orch) -- (rl0);

  \node[sbox, fill=green!15] (s1) at (-1.1, 0.6) {Subagent A};
  \node[sbox, fill=green!15] (s2) at ( 1.1, 0.6) {Subagent B};
  \draw[-{Stealth}] (rl0) -- (s1);
  \draw[-{Stealth}] (rl0) -- (s2);

  \node[sbox, fill=orange!15, minimum width=1.3cm] (rl2) at (-1.1,-0.15) {\scriptsize ReAct};
  \node[sbox, fill=orange!15, minimum width=1.3cm] (rl3) at ( 1.1,-0.15) {\scriptsize ReAct};
  \draw[-{Stealth}] (s1) -- (rl2);
  \draw[-{Stealth}] (s2) -- (rl3);

  \node[font=\scriptsize\bfseries] at (0, 2.7) {Multi-agent};
\end{tikzpicture}
\end{column}
\end{columns}

\end{frame}

